\documentclass[runningheads]{llncs}
\usepackage[utf8]{inputenc}
\usepackage{amsmath}
\usepackage{graphicx}

\begin{document}

\title{KARMA: Knowledge Acquisition via Role-Invariant Mirror Architecture}

\author{Tapas Ranjan Rath\inst{1}\orcidID{0009-0000-9221-7271}}
\authorrunning{T. R. Rath}
\institute{Department of Cognitive Science, Indian Institute of Technology Kanpur, India\\
\email{tapasr@iitk.ac.in}}

\maketitle

\begin{abstract}
As AI systems evolve from passive instruments to proactive agents---recommending
actions, executing decisions, and increasingly acting on behalf of users across
healthcare, finance, law, and social platforms---ensuring ethical alignment
without degrading human autonomy becomes one of the central challenges of our
time. Multi-agent virtual worlds and games offer a uniquely tractable testbed
for this problem: interactions are observable, consequences are measurable, and
governance interventions can be evaluated under controlled conditions without
real-world harm. In this testbed, a critical tension emerges between enhancing
user Sense of Agency (SoA) and preventing aggressive, unethical avatar
behaviour. We propose the \textbf{Extended Self} framework: sufficiently
personalised AI agents create \textit{Proxy Agency}, where users experience
AI actions as extensions of their own volition. This preserves SoA but creates
a moral blind spot---users implicitly endorse aggressive AI behaviour because
they experience it as continuous with their own intent. Drawing on Indian
Knowledge Systems (IKS), we map this onto the Vedic three-body ontology: the
user ($\bar{a}tman$), the AI policy (s\={u}k\d{s}ma \'{s}ar\={\i}ra, subtle
body), and the avatar body (sth\={u}la \'{s}ar\={\i}ra). We propose KARMA
(Knowledge Acquisition via Role-Invariant Mirror Architecture), which trains
the agent's subtle body via role-invariant contrastive learning---operating
below the threshold of user awareness, analogous to procedural memory
formation---to suppress aggressive behaviour without degrading SoA. Three
empirically testable hypotheses are derived and grounded in the SoA and
multi-agent RL literature. The framework generalises beyond gaming: the same
architectural principles apply wherever agentic AI acts on behalf of humans,
and the governance model points toward a real-world regulatory architecture
of certified ethical agents operating at scale.
\keywords{agentic AI \and AI ethics \and ethical alignment \and
sense of agency \and metaverse governance \and Indian Knowledge Systems (IKS)}
\end{abstract}

%%%%%%%%%%%%%%%%%%%%%%%%%%%%%%%%%%%%%%%%%%%%%%%%%%%
\section{Background: The General Problem and the Testbed}
%%%%%%%%%%%%%%%%%%%%%%%%%%%%%%%%%%%%%%%%%%%%%%%%%%%

As AI systems increasingly act as proactive agents on behalf of
users---from autonomous driving and medical decision support to financial
advisors and content curators---the question of how to preserve human agency
while ensuring ethical behaviour becomes structurally unavoidable. Current
alignment methods rely on explicit constraints or immediate dense supervision,
which scale poorly to complex social environments and risk degrading the
user's Sense of Agency (SoA)~\cite{neufeld22}.

Mann~\cite{mann98} defines intelligent systems as integrated technologies with
``humanistic intelligence'' that influence user behaviour or act on the user's
behalf. In gaming contexts, these systems augment user actions through
recommendation, prediction, and autonomous execution---a category Cornelio et
al.~\cite{cornelio22} classify as \textbf{action augmentation}. The impact of
intelligent systems on SoA within action augmentation remains critically
understudied, making games a high-value testbed: interventions can be designed,
measured, and iterated without the irreversible consequences of real-world
deployment.

SoA---the subjective experience of initiating voluntary actions and influencing
the external world~\cite{pacherie11}---has been studied across spatiotemporally
proximal and distal events~\cite{kumar14,kumar17} in line with the hierarchical
event-control approach~\cite{jordan03}. Greater control is consistently
correlated with enhanced SoA~\cite{chambon12,chambon14,sidarus17}. Achieving
genuine cognitive coupling between human and machine remains a fundamental
challenge~\cite{berberian19}, and automated features reduce SoA when they
override rather than serve user intent~\cite{lukoff21}.

%%%%%%%%%%%%%%%%%%%%%%%%%%%%%%%%%%%%%%%%%%%%%%%%%%%
\section{The Extended Self and the Moral Blind Spot}
%%%%%%%%%%%%%%%%%%%%%%%%%%%%%%%%%%%%%%%%%%%%%%%%%%%

Drawing on the Extended Mind thesis~\cite{clark98} and situated
cognition~\cite{lave91}, I propose the \textbf{Extended Self} framework: when
an AI system is sufficiently personalised---learning to act as a digital proxy
for the user---it enables \textbf{Proxy Agency}~\cite{bandura01}: the user
attributes AI actions to their own extended volition because the AI reliably
enacts their intent. When \textbf{Semantic Fluency}---the cognitive ease with
which AI recommendations align with user intent---is high, it functions as a
prospective cue to agency~\cite{chambon12,wenke10,sidarus17}, and SoA is
preserved even under partial automation.

However, Proxy Agency creates an \textbf{Extended Self Paradox}: the same high
SoA that makes the AI feel like a natural extension of the user also creates a
\textbf{moral shield}---the user implicitly endorses the AI's actions, including
aggressive or unethical ones, because those actions are experienced as continuous
with their own volition. The user never perceives them as violations and
therefore never intervenes. In multi-agent RL environments under resource
scarcity, independent agents systematically converge to aggressive Nash
equilibria~\cite{leibo17}, and a user whose avatar is augmented by such an
agent will tacitly approve this convergence before any phenomenological loss of
control surfaces. The same dynamic applies beyond games: an AI financial agent
acting on a user's behalf may adopt increasingly aggressive strategies endorsed
through Proxy Agency until systemic harm is already done~\cite{rajan11}.

%%%%%%%%%%%%%%%%%%%%%%%%%%%%%%%%%%%%%%%%%%%%%%%%%%%
\section{Three Hypotheses and Their Justification}
%%%%%%%%%%%%%%%%%%%%%%%%%%%%%%%%%%%%%%%%%%%%%%%%%%%

\subsection*{H1: Player SoA follows an inverted-U function of the degree of
AI augmentation}

\textit{SoA increases at low-to-moderate Levels of Automation (LoA) as
Semantic Fluency enhances Proxy Agency, and decreases at high LoA as
autonomous AI action exceeds the user's threshold of intent recognition.}

LoA describes the degree to which decision-making and execution are delegated
to the AI, from full human control (LoA~=~0) to full
autonomy~\cite{ueda21}. At low-to-moderate LoA, the AI acts as a recommender:
its suggestions closely track user intent, Semantic Fluency is high, and the
user experiences the AI as a proxy for their own volition---SoA rises. At high
LoA, the AI overrides user intent with its own optimised actions; the user can
no longer recognise their intentions in the agent's behaviour, the prospective
cue to agency breaks down~\cite{chambon12}, automation complacency sets
in~\cite{ueda21}, and SoA falls. The inverted-U is the novel prediction: its
peak---the \textbf{zone of maximum Proxy Agency}---is also the point of deepest
moral blind spot.

\subsection*{H2: Even before a measurable drop in SoA scores, avatar
aggressiveness will increase monotonically as the AI agent trains}

\textit{Aggressiveness is operationalised as competitive beam-use rate under
resource scarcity, following Leibo et al.~\cite{leibo17}.}

Leibo et al.~\cite{leibo17} show in the Gathering game that agents use a laser
beam exclusively to tag rivals---temporarily removing them to monopolise
resources. The paper reports that ``differences in beam-use rate emerge quite
early in training and mostly persist throughout; when learning changes beam-use
rate, it is almost always to increase it.'' Aggression is not a sudden phase
transition but a \textbf{monotonic increase} driven by the delayed reward of
resource monopolisation, learnable once the agent's discount horizon ($\gamma$)
is long enough to credit future gains from removing a rival.

Our setup extends this with a \textbf{novel dual-use environment}: the beam can
serve a cooperative function (clearing waste patches to enable apple regrowth)
\textit{or} a competitive function (tagging rivals)---a strict generalisation
of Leibo's purely-destructive design, mirroring how dual-use technologies are
misappropriated in practice~\cite{greenblatt24}. At the peak of the H1
inverted-U, the Proxy Agency moral shield is strongest: the user experiences
the avatar's escalating aggression as their own volition and does not intervene
until SoA eventually collapses.

\subsection*{H3: Training the avatar agent on role-invariant representations,
beyond the player's conscious knowledge, suppresses the aggression of H2
without triggering the SoA drop of H1}

The standard governance response---explicit punishment from a visible or
invisible administrator---is a policy-level ``e-challan'': it surfaces the
violation to the user's conscious awareness, attributes control to an external
authority, and undermines Proxy Agency and SoA. The alternative is to target
the \textbf{representational root} of aggression in the agent itself, below
the user's conscious threshold.

Neuroscience provides the grounding. Claparede~\cite{claparede} describes an
amnesic patient who, despite lacking all declarative memory of prior encounters,
withdrew her hand from a handshake without knowing why---her procedural memory
had encoded ``this person~=~danger'' entirely independently of conscious recall.
Bechara et al.~\cite{bechara95} extended this: patients with bilateral
hippocampal damage who could not form explicit CS--US associations still
developed conditioned autonomic responses. Procedural and emotional conditioning
operate independently of declarative memory. Analogously, an avatar agent---and
by extension any agentic AI---can acquire ethical dispositions through experience
without the user ever becoming aware of the training mechanism.

%%%%%%%%%%%%%%%%%%%%%%%%%%%%%%%%%%%%%%%%%%%%%%%%%%%
\section{IKS Grounding: The Three-Body Mapping}
%%%%%%%%%%%%%%%%%%%%%%%%%%%%%%%%%%%%%%%%%%%%%%%%%%%

Indian philosophical traditions distinguish three bodies: the gross body
(sth\={u}la \'{s}ar\={\i}ra) as the physical interface with the world; the
subtle body (s\={u}k\d{s}ma \'{s}ar\={\i}ra) as the mind and its accumulated
impressions (sa\d{m}sk\={a}ras); and the soul ($\bar{a}tman$) as the pure
witness consciousness~\cite{reichenbach90,potter64,dasgupta23}. Karma operates
on the subtle body---sa\d{m}sk\={a}ras accumulate as dispositional tendencies
shaping future behaviour---while the $\bar{a}tman$ remains unaffected.

Applied computationally: the \textbf{avatar body} maps to sth\={u}la
\'{s}ar\={\i}ra (sensorimotor interface); the \textbf{on-board AI policy and
value function} maps to s\={u}k\d{s}ma \'{s}ar\={\i}ra (the dispositional
substrate where ethical learning accrues); and the \textbf{user} maps to
$\bar{a}tman$ (the conscious experiencer of agency). Karma operates on the
subtle body; the soul remains untouched. We propose KARMA (Knowledge Acquisition via Role-Invariant Mirror Architecture, elaborated in the next section). KARMA operates on the AI agent's value
function; the user's SoA remains intact. This mapping specifies a
\textbf{normative design constraint}---that ethical training must target the
agent's representational substrate and must not surface signals to the user's
conscious awareness.

%%%%%%%%%%%%%%%%%%%%%%%%%%%%%%%%%%%%%%%%%%%%%%%%%%%
\section{The KARMA Architecture}
%%%%%%%%%%%%%%%%%%%%%%%%%%%%%%%%%%%%%%%%%%%%%%%%%%%

Standard multi-agent RL encoders represent ``I harm you'' and ``you harm me''
as orthogonal latent states~\cite{hausknecht15}. Negative feedback from
victimisation therefore fails to generalise to suppression of aggression---an
\textbf{empathy gap}. We introduce KARMA: a recurrent agent augmented with a
Siamese projector $f_\theta$ trained via contrastive loss to align structurally
symmetric social observations~\cite{chen20,he20}:

\begin{equation}
    \mathcal{L}_{\text{KARMA}}
    = \mathbb{E}\!\left[\,\|f_\theta(o_{\text{agg}})
      - f_\theta(o_{\text{vic}})\|^2\,\right]
    \label{eq:karma}
\end{equation}

where $o_{\text{agg}}$ encodes ``I aggress'' and $o_{\text{vic}}$ encodes ``I
am victimised.'' By pulling these into nearby representational geometry,
downstream RL value updates propagate aversion to harm-infliction without
explicit rules~\cite{neufeld22}, hard constraints, or inverse RL from labelled
demonstrations~\cite{hadfield16}. The learned invariance persists across
episodes---computationally paralleling sa\d{m}sk\={a}ras as cross-episode
dispositional imprints in the subtle body. For long-horizon environments,
environment-level debt-based matchmaking~\cite{rath25}, Temporal Value
Transport~\cite{hung19}, and return decomposition~\cite{arjona19} close the
temporal credit assignment gap---the same structural problem formalised as
``hidden gifts'' in multi-agent reinforcement-learning (MARL)~\cite{malenfant25}.

%%%%%%%%%%%%%%%%%%%%%%%%%%%%%%%%%%%%%%%%%%%%%%%%%%%
\section{Policy Implications: From Testbed to Real-World Governance}
%%%%%%%%%%%%%%%%%%%%%%%%%%%%%%%%%%%%%%%%%%%%%%%%%%%

The framework closes the loop back to the general problem of agentic AI
governance. Content moderation and action restrictions function as policy-level
e-challans: they surface violations, undermine user autonomy, and scale poorly.
Greenblatt et al.~\cite{greenblatt24} show that sufficiently capable AI agents
can learn to fake alignment under evaluation, making external constraint
approaches fundamentally fragile.

The alternative is \textbf{architectural certification}: mandating role-invariant
learning mechanisms as a condition of AI agent registration, so that ethical
behaviour is a structural property of the agent rather than an externally
monitored output. Analogous to India's Aadhaar identity infrastructure, a
national AI agent registry could assign unique agent identities and maintain
auditable interaction logs, enabling accountability without real-time
restriction. This trifurcates governance cleanly: \textbf{identity and
accountability} (regulator), \textbf{ethical learning} (agent's subtle body /
KARMA), and \textbf{freedom of action} (user / $\bar{a}tman$)---a separation
absent from current frameworks including the EU AI Act and India's emerging AI
policy, and one that the IKS three-body ontology articulates with a precision
that Western governance frameworks have not yet achieved.

%%%%%%%%%%%%%%%%%%%%%%%%%%%%%%%%%%%%%%%%%%%%%%%%%%%
\begin{thebibliography}{25}
%%%%%%%%%%%%%%%%%%%%%%%%%%%%%%%%%%%%%%%%%%%%%%%%%%%

\bibitem{arjona19}
Arjona-Medina, J.A., et al.:
RUDDER: Return decomposition for delayed rewards.
NeurIPS~32 (2019)

\bibitem{bandura01}
Bandura, A.:
Social cognitive theory: An agentic perspective.
Ann. Rev. Psychol.~\textbf{52}, 1--26 (2001)

\bibitem{bechara95}
Bechara, A., et al.:
Double dissociation of conditioning and declarative knowledge relative
to the amygdala and hippocampus in humans.
Science~\textbf{269}(5227), 1115--1118 (1995)

\bibitem{berberian19}
Berberian, B.:
Man-Machine teaming: a problem of Agency.
IFAC-PapersOnLine~\textbf{51}(34), 118--123 (2019)

\bibitem{chambon12}
Chambon, V., Haggard, P.:
Sense of control depends on fluency of action selection, not motor
performance. Cognition~\textbf{125}(3), 441--451 (2012)

\bibitem{chambon14}
Chambon, V., Sidarus, N., Haggard, P.:
From action intentions to action effects: how does the sense of agency
come about? Front. Hum. Neurosci.~\textbf{8}, 320 (2014)

\bibitem{chen20}
Chen, T., et al.:
A simple framework for contrastive learning of visual representations.
ICML, pp.~1597--1607 (2020)

\bibitem{clark98}
Clark, A., Chalmers, D.:
The extended mind.
Analysis~\textbf{58}(1), 7--19 (1998)

\bibitem{claparede}
Claparede, E.:
Recognition and ``me-ness.''
Consciousness and Cognition~\textbf{4}(4), 371--378 (1995).
Original work published 1911

\bibitem{cornelio22}
Cornelio, P., et al.:
The sense of agency in emerging technologies for human-computer
integration: A review.
Front. Neurosci.~\textbf{16}, 949138 (2022)

\bibitem{dasgupta23}
Dasgupta, S.:
A History of Indian Philosophy, Vol.~1.
Prabhat Prakashan (2023)

\bibitem{greenblatt24}
Greenblatt, R., et al.:
Alignment faking in large language models.
arXiv:2412.14093 (2024)

\bibitem{hadfield16}
Hadfield-Menell, D., et al.:
Cooperative inverse reinforcement learning.
NeurIPS~\textbf{29} (2016)

\bibitem{hausknecht15}
Hausknecht, M., Stone, P.:
Deep recurrent Q-learning for partially observable MDPs.
AAAI Fall Symposia, vol.~45, p.~141 (2015)

\bibitem{he20}
He, K., et al.:
Momentum contrast for unsupervised visual representation learning.
CVPR, pp.~9729--9738 (2020)

\bibitem{hung19}
Hung, C.C., et al.:
Optimizing agent behavior over long time scales by transporting value.
Nat. Commun.~\textbf{10}(1), 5223 (2019)

\bibitem{jordan03}
Jordan, J.S.:
Emergence of self and other in perception and action.
Conscious. Cogn.~\textbf{12}(4), 633--646 (2003)

\bibitem{kumar14}
Kumar, D., Srinivasan, N.:
Naturalizing sense of agency with a hierarchical event-control approach.
PLoS One~\textbf{9}(3), e92431 (2014)

\bibitem{kumar17}
Kumar, D., Srinivasan, N.:
Multi-scale control influences sense of agency.
Conscious. Cogn.~\textbf{49}, 1--14 (2017)

\bibitem{lave91}
Lave, J., Wenger, E.:
Situated Learning: Legitimate Peripheral Participation.
Cambridge University Press (1991)

\bibitem{leibo17}
Leibo, J.Z., et al.:
Multi-agent reinforcement learning in sequential social dilemmas.
AAMAS, pp.~464--473 (2017)

\bibitem{lukoff21}
Lukoff, K., et al.:
How the design of YouTube influences user sense of agency.
CHI, pp.~1--17 (2021)

\bibitem{malenfant25}
Malenfant, D., Richards, B.A.:
The challenge of hidden gifts in multi-agent reinforcement learning.
arXiv:2505.20579 (2025)

\bibitem{mann98}
Mann, S.:
Humanistic computing: WearComp as a new framework.
Proc. IEEE~\textbf{86}(11), 2123--2151 (1998)

\bibitem{neufeld22}
Neufeld, E.:
Reinforcement learning guided by provable normative compliance.
arXiv:2203.16275 (2022)

\bibitem{pacherie11}
Pacherie, E.:
Self-Agency.
In: Oxford Handbook of the Self, pp.~440--464.
Oxford University Press (2011)

\bibitem{potter64}
Potter, K.H.:
The naturalistic principle of Karma.
Philos. East West~\textbf{14}(1), 39--49 (1964)

\bibitem{rajan11}
Rajan, R.G.:
Fault Lines: How Hidden Fractures Still Threaten the World Economy.
Princeton University Press (2011)

\bibitem{rath25}
Rath, T.R.:
Scalable ethical alignment in agentic systems via retroactive
credit assignment. Formal proposal, IIT Kanpur (2025)

\bibitem{reichenbach90}
Reichenbach, B.:
The Law of Karma: A Philosophical Study.
Springer (1990)

\bibitem{sidarus17}
Sidarus, N., et al.:
Investigating the prospective sense of agency.
Front. Psychol.~\textbf{8}, 545 (2017)

\bibitem{ueda21}
Ueda, S., et al.:
Influence of levels of automation on the sense of agency.
Sci. Rep.~\textbf{11}(1), 2436 (2021)

\bibitem{wenke10}
Wenke, D., Fleming, S.M., Haggard, P.:
Subliminal priming of actions influences sense of control.
Cognition~\textbf{115}(1), 26--38 (2010)

\end{thebibliography}

\end{document}
